\documentclass[10pt,a4paper]{amsart}
\usepackage[margin=19mm]{geometry}
\usepackage{amsmath,amssymb,mathtools}
\usepackage[scr=rsfs,cal=euler]{mathalfa}
\usepackage{xfrac}
\usepackage[unicode,hidelinks,bookmarks=false]{hyperref}
\newcommand{\ie}{i.e.\ }
\newcommand{\cf}{cf.\ }
\newcommand{\resp}{respectively\ }
\newcommand{\Subsection}{\subsection}
\providecommand{\nobreakdash}{\nobreak\hbox{-}}
\setlength{\emergencystretch}{2em}
\multlinegap0pt
\usepackage{amssymb}
\usepackage{enumerate}
\newtheorem{theorem}{Theorem}
\title{\LARGE \bf Far field refraction problem with loss of energy in negative refractive index material}
\author{Haokun Sui, Feida Jiang}
\date{Source: arXiv:2603.11092v1 (CC BY 4.0)}
\begin{document}
\maketitle

\noindent\emph{Source and licence.} \LARGE \bf Far field refraction problem with loss of energy in negative refractive index material. Version arXiv:2603.11092v1 (CC BY 4.0), distributed by arXiv under the Creative Commons Attribution 4.0 International licence, http://creativecommons.org/licenses/by/4.0/, as recorded in the arXiv licence metadata for this exact version and shown on the abstract page https://arxiv.org/abs/2603.11092v1. Full licence notice: \texttt{licenses/arxiv-2603.11092-CC-BY-4.0.txt}.\par\smallskip

\noindent\textit{Editorial scope.} Counted unit: the article's theorem 5.6 (source0.tex lines 1400--1407). The article's complete original proof of it is delivered with it (source0.tex lines 1409--1597, one \texttt{proof} environment of 189 lines). No mathematics is written, repaired or re-derived: the statement and the proof are the article's own lines, in the article's own notation, and no retained line is substituted or re-flowed.  Retained uncounted as the source-local context the proof needs: lines 158--160; lines 1347--1349; lines 1365--1367; lines 1394--1396.  Nothing was omitted from any retained span, so the package carries no omission record.  No retained line contains a citation, so the wrapper carries no bibliography and no bibliography entry.  No retained line contains a figure environment, an image-inclusion command or drawing code, so no image is delivered and the package declares no asset dependency.  Editorial formatting only, in addition: an AMS \texttt{amsart} preamble that restates the article's own declarations and macros used by the retained lines, namely the environments \texttt{theorem} on separate counters, together with the standard packages those lines need.  The retained environments print in this document as Theorem 1 in order of appearance, whereas the article prints the same items with the numbering of its own full document.  The complete native source of the article as distributed in the arXiv e-print for this exact version is bundled unchanged as \texttt{sources/196192/source0.tex}.
\begin{equation}\label{2.5}
  \Phi(t) = t + \vert \kappa \vert \sqrt{1 - \kappa^{-2}(1 - t^{2})},
\end{equation}

\begin{equation}\label{5.1}
  Y = \frac{1}{\kappa}(X - \Phi(x \cdot \nu)\nu),
\end{equation}

\begin{equation}\label{5.2}
  \det J = \frac{1}{y_{n}} \det Dy,
\end{equation}

\begin{equation}\label{5.5}
   \vert \det J \vert = \frac{dS_{\bar{\Omega}^{\ast}}}{dS_{\mathcal{V}}} \leq \frac{f(x)t_{\Gamma}(x)}{\sqrt{1 - \vert x \vert^{2}}g(T(x))}.
\end{equation}

\begin{theorem}
  Suppose a refractor $\Gamma$ is defined by $\rho$, and $\rho$ is the weak solution to the refractor problem in negative refractive index material with loss of energy with emitting illumination intensity $f \in L^{1}(\bar{\Omega})$ and prescribed refracted illumination intensity $g \in L^{1}(\bar{\Omega}^{\ast})$, then we have
  \begin{equation}\label{5.6}
    \vert \det (D^{2}\rho + C^{-1}B) \vert \leq \frac{f(x)t_{\Gamma}(x)\vert \kappa \vert^{n - 2} \omega}{g(T(x))h^{n - 1}\left(1 - h^{-1}\left(\dfrac{\rho}{1 - \vert x \vert^{2}}x - D\rho\right)\cdot D_{p}h\right)},
  \end{equation}
  where $C^{-1}$ is given in (\ref{5.25}), $B$ is given by (\ref{5.23}), $h$ and $\omega$ are defined in (\ref{5.14}) and (\ref{5.15}) correspondingly.
  \label{thm5.1}
\end{theorem}

\begin{proof}
  In order to prove (\ref{5.6}), we first need to derive the unit outer normal $\nu$ to the surface $\Gamma$ at $\rho(X)X$.

  Write $\nu = (\nu',\nu_{n})$, for $\partial_{x_{k}}((x,x_{n})\rho(x))$ are tangential to the graph of the refractor $\Gamma$ for $k = 1,2,\ldots,n-1$, then we have
  $$\partial_{x_{k}}((x,x_{n})\rho(x))\cdot \nu = 0$$
  for $k = 1,2,\ldots,n-1$. Hence we have
  \begin{equation}\label{5.7}
    \rho \sum_{i = 1}^{n-1}\delta_{ik}\nu_{i} + \partial_{x_{k}}\rho \sum_{i = 1}^{n-1}x_{i}\nu_{i} = \left(\rho \frac{x_{k}}{\sqrt{1 - \vert x \vert^{2}}} - \sqrt{1 - \vert x \vert^{2}}\partial_{x_{k}}\rho\right)
  \end{equation}
  for $k = 1,2,\ldots,n-1$. Using the tensor product, (\ref{5.7}) can be written in matrix form
  \begin{equation}\label{5.8}
    (\rho I + D\rho \otimes x)(\nu')^{T} = \left(\rho \frac{x^{T}}{\sqrt{1 - \vert x \vert^{2}}} - \sqrt{1 - \vert x \vert^{2}}(D\rho)^{T}\right)\nu_{n}.
  \end{equation}
  According to Shermann-Morrison formula, we have
  \begin{equation*}
    (\rho I + D\rho \otimes x)^{-1} = \rho^{-1}\left(I - \frac{D\rho \otimes x}{\rho + D\rho \cdot x}\right).
  \end{equation*}
  Notice that for any row vectors $a,b$ and $c$, we have $(a \otimes b)c^{T} = (b \cdot c)a^{T}$, then we have
  \begin{align*}
    (\nu')^{T} & = \rho^{-1}\left(I - \frac{D\rho \otimes x}{\rho + D\rho \cdot x}\right)\left(\rho \frac{x^{T}}{\sqrt{1 - \vert x \vert^{2}}} - \sqrt{1 - \vert x \vert^{2}}(D\rho)^{T}\right)\nu_{n} \\
     & = \rho^{-1}\left(\frac{\rho}{\sqrt{1 - \vert x \vert^{2}}}x^{T} - \sqrt{1 - \vert x \vert^{2}}(D\rho)^{T} - \frac{\rho}{\sqrt{1 - \vert x \vert^{2}}(\rho + D\rho \cdot x)}(D\rho \otimes x)x^{T}\right. \\
     & \quad \left. + \frac{\sqrt{1 - \vert x \vert^{2}}}{\rho + D\rho \cdot x}(D\rho \otimes x)(D\rho)^{T}\right)\nu_{n} \\
     & = \rho^{-1}\left(\frac{\rho}{\sqrt{1 - \vert x \vert^{2}}}x^{T} - (\sqrt{1 - \vert x \vert^{2}} + \frac{\vert x \vert^{2}\rho}{\sqrt{1 - \vert x \vert^{2}}(\rho + D\rho \cdot x)}\right. \\
     & \quad \left. - \frac{\sqrt{1 - \vert x \vert^{2}}}{\rho + D\rho \cdot x}(x \cdot D\rho))(D\rho)^{T}\right)\nu_{n} \\
     & = \frac{1}{\sqrt{1 - \vert x \vert^{2}}}\left(x^{T} - \frac{1}{\rho + D\rho \cdot x}(D\rho)^{T}\right)\nu_{n}.
  \end{align*}
  Hence we have
  \begin{equation}\label{5.9}
    \nu = \left(\frac{1}{\sqrt{1 - \vert x \vert^{2}}}\left(x - \frac{1}{\rho + D\rho \cdot x}D\rho\right),1\right)\nu_{n}.
  \end{equation}
  From (\ref{5.9}), we have
  \begin{equation}\label{5.10}
    X\cdot \nu = \frac{1}{\sqrt{1 - \vert x \vert^{2}}}\left(\frac{\rho}{\rho + D\rho \cdot x}\right)\nu_{n}.
  \end{equation}
  Since $\nu$ is unit outer normal to $\Gamma$ at $\rho(X)X$, then we have $X\cdot \nu \geq 0$ and $\vert \nu' \vert^{2} + \nu_{n}^{2} = 1$, then from (\ref{5.9}), we have
  \begin{equation}\label{5.11}
    \left(\frac{\rho^{2} - (x \cdot D\rho)^{2} + \vert D\rho \vert^{2}}{(1 - \vert x \vert^{2})(\rho + D\rho \cdot x)^{2}}\right)\nu_{n}^{2} = 1.
  \end{equation}
  According to (\ref{5.11}), we obtain
  \begin{equation}\label{5.12}
    \nu_{n} = \pm \vert \rho + D\rho \cdot x \vert\sqrt{\frac{1-\vert x \vert^{2}}{\rho^{2} - (x \cdot D\rho)^{2} + \vert D\rho \vert^{2}}}.
  \end{equation}
  Then from (\ref{5.9}), we get
  \begin{equation}\label{5.13}
    \nu = \pm \frac{\vert \rho + D\rho \cdot x \vert}{\rho + D\rho \cdot x}\left(\frac{-D\rho +\rho + D\rho \cdot x}{\sqrt{\rho^{2} - (x \cdot D\rho)^{2} + \vert D\rho \vert^{2}}}x,\frac{\sqrt{1 - \vert x \vert^{2}}(\rho + D\rho \cdot x)}{\sqrt{\rho^{2} - (x \cdot D\rho)^{2} + \vert D\rho \vert^{2}}}\right).
  \end{equation}
  Besides, from (\ref{5.13}), (\ref{5.10}) can be written as
  \begin{equation*}
    X\cdot \nu = \pm \frac{\vert \rho + D\rho \cdot x \vert}{\rho + D\rho \cdot x}\frac{\rho}{\sqrt{\rho^{2} - (x \cdot D\rho)^{2} + \vert D\rho \vert^{2}}}.
  \end{equation*}

  Now we can derive the inequality involving a Monge-Amp\`ere type operator satisfied by $\rho$. For simplicity, we introduce two functions:
  \begin{equation}\label{5.14}
    h(x,z,p) = \frac{\Phi\left(\dfrac{z}{\sqrt{z^{2}+\vert p \vert^{2}-(p \cdot x)^{2}}}\right)}{\sqrt{z^{2}+\vert p \vert^{2}-(p \cdot x)^{2}}}
  \end{equation}
  and
  \begin{equation}\label{5.15}
   \omega(x,z,p) = 1 - h(x,z,p)(z + p \cdot x),
  \end{equation}
  where $\Phi$ is defined in (\ref{2.5}). Then from (\ref{5.1}), we have
  \begin{equation}\label{5.16}
    y_{i} = \frac{1}{\kappa}\left[\omega(x,\rho(x),D\rho(x))x_{i} + h(x,\rho(x),D\rho(x))\rho_{x_{i}}\right], \quad 1 \leq i \leq n-1,
  \end{equation}
  and
  \begin{equation}\label{5.17}
    y_{n} = \frac{1}{\kappa} \omega(x,\rho(x),D\rho(x)) \sqrt{1 - \vert x \vert^{2}}.
  \end{equation}
  For $1 \leq i,j \leq n-1$, differentiating $y_{i}$ with respect to $x_{j}$, we have
  \begin{equation}\label{5.18}
    \begin{aligned}
    \partial_{j}y_{i} & = \frac{1}{\kappa}\left[\omega \delta_{ij} + x_{i}(\omega_{x_{j}} + \omega_{z} \rho_{x_{j}} + \sum_{k = 1}^{n-1}\omega_{p_{k}} \rho_{x_{k}x_{j}})\right.  
      \left. + h \rho_{x_{i}x_{j}} + \rho_{x_{i}}(h_{x_{j}} + h_{z} \rho_{x_{j}} + \sum_{k=1}^{n-1} h_{p_{k}} \rho_{x_{k}x_{j}})\right].
    \end{aligned}
  \end{equation}
  For $x, D\rho, D_{x}\omega, D_{p}\omega, D_{z}\omega, D_{z}\omega, D_{x}h$ and $D_{p}h$ are row vectors, then (\ref{5.18}) can be written in matrix form
  \begin{equation}\label{5.19}
    \begin{aligned}
    Dy & = \frac{1}{\kappa}[\omega I + x \otimes D_{x}\omega + \omega_{z}x \otimes D\rho + D\rho \otimes D_{x}h  \\
    & \quad + h_{z}D\rho \otimes D\rho + (x \otimes D_{p}\omega)D^{2}\rho + hD^{2}\rho + (D\rho \otimes D_{p}h)D^{2}\rho].
    \end{aligned}
  \end{equation}
  Let
  \begin{equation}\label{5.20}
    B(x) = \omega I + x \otimes D_{x}\omega + \omega_{z}x \otimes D\rho + D\rho \otimes D_{x}h + h_{z}D\rho \otimes D\rho,
  \end{equation}
  and
  \begin{equation}\label{5.21}
    C(x) = x \otimes D_{p}\omega + hI + D\rho \otimes D_{p}h.
  \end{equation}
  Then (\ref{5.19}) can be written as
  \begin{equation}\label{5.22}
    Dy = \frac{1}{\kappa}[B(x) + C(x)D^{2}\rho].
  \end{equation}
  From (\ref{5.15}), we have $D_{x}\omega = -D_{x}h -hp$, $\omega_{z} = -h_{z}(z + p \cdot x) - h$ and $D_{p}\omega = -D_{p}h(z + p \cdot x) - hx$, then we have
  \begin{equation}\label{5.23}
    \begin{aligned}
    B(x) & = [1 - (\rho + D\rho \cdot x)h]I - [(\rho + D\rho \cdot x)x - D\rho]\otimes D_{x}h  \\
    & \quad - x\otimes [(2h + h_{z}((\rho + D\rho \cdot x)))D\rho] + h_{z}D\rho \otimes D\rho,
    \end{aligned}
  \end{equation}
  and
  \begin{equation}\label{5.24}
    \begin{aligned}
    C(x) & = [(-(\rho + D\rho \cdot x)x + D\rho) \otimes D_{p}h] - h[(x \otimes x) - I] \\
    & = h[(h^{-1}(-(\rho + D\rho \cdot x)x + D\rho)\otimes D_{p}h) + (((-x) \otimes x) + I)] \\
    & = h(M_{1} + M_{2}) \\
    & = h M_{2}(I + M_{2}^{-1}M_{1}),
    \end{aligned}
  \end{equation}
  where
  \begin{equation*}
    M_{1} = -h^{-1}((\rho + D\rho \cdot x)x - D\rho) \otimes D_{p}h \quad \text{and} \quad M_{2} = ((-x)\otimes x) + I.
  \end{equation*}
  From Shermann-Morrison formula, we have
 $
    M_{2}^{-1} = I + \frac{x \otimes x}{1 - \vert x \vert^{2}},
 $
  and
 $
    C^{-1} = \frac{1}{h}(I + M_{2}^{-1}M_{1})M_{2}^{-1}.
  $
  We may assume that $v = h^{-1}[(\rho + D\rho \cdot x)x - D\rho]$, then $I + M_{2}^{-1}M_{1} = I + (-M_{2}^{-1}v^{T})D_{p}h$, then we have
  \begin{equation*}
    N := (I + M_{2}^{-1}M_{1})^{-1} = I + \frac{(-M_{2}^{-1}v^{T})D_{p}h}{1 - (-M_{2}^{-1}v^{T})^{T}\cdot D_{p}h},
  \end{equation*}
  hence
  \begin{equation}\label{5.25}
    C^{-1} = \frac{1}{h}N\left(I + \frac{x \otimes x}{1 - \vert x \vert^{2}}\right).
  \end{equation}
  Now we calculate the matrix $N$ accurately. We have
  \begin{align*}
    M_{2}^{-1}v^{T} & = v^{T} + \frac{1}{1 - \vert x \vert^{2}} x^{T}xv^{T} \\
      & = h^{-1}\left[(\rho + D\rho \cdot x)x^{T} - (D\rho)^{T} + \frac{1}{1 - \vert x \vert^{2}}(\vert x \vert^{2}(\rho + D\rho \cdot x)x^{T} - (D\rho \cdot x)x^{T})\right] \\
      & = h^{-1}\left[\frac{\rho}{1 - \vert x \vert^{2}}x^{T} - (D\rho)^{T}\right].
  \end{align*}
  Hence, we have
  \begin{equation}\label{5.26}
    N = I + \frac{h^{-1}\left(\dfrac{\rho}{1 - \vert x \vert^{2}}x - D\rho\right)\otimes D_{p}h}{1 - h^{-1}\left(\dfrac{\rho}{1 - \vert x \vert^{2}}x - D\rho\right)\cdot D_{p}h},
  \end{equation}
  again from Shermann-Morrison formula, we obtain
    \begin{equation}\label{5.27}
    \det N = \frac{1}{1 - h^{-1}\left(\dfrac{\rho}{1 - \vert x \vert^{2}}x - D\rho\right)\cdot D_{p}h}.
  \end{equation}
  Combining (\ref{5.25}) with (\ref{5.27}), we have
    \begin{equation}\label{5.28}
    \det C = \frac{1}{\det C^{-1}} = h^{n-1}(1 - \vert x \vert^{2})\left[1 - h^{-1}\left(\frac{\rho}{1 - \vert x \vert^{2}}x - D\rho\right)\cdot D_{p}h\right].
  \end{equation}
  Substituting (\ref{5.26}) into (\ref{5.25}), we have
    \begin{align*}
    C^{-1} & = \frac{1}{h}\left[I + \frac{x\otimes x}{1 - \vert x \vert^{2}} + \frac{(M_{2}^{-1}v^{T})D_{p}h}{1 - (M_{2}^{-1}v^{T})^{T}\cdot D_{p}h} + \left(\frac{(M_{2}^{-1}v^{T})D_{p}h}{1 - (M_{2}^{-1}v^{T})^{T}\cdot D_{p}h}\right)\left(\frac{x \otimes x}{1 - \vert x \vert^{2}}\right)\right] \\
      & = \frac{1}{h}\left[I + \frac{x\otimes x}{1 - \vert x \vert^{2}} + \frac{1}{1 - (M_{2}^{-1}v^{T})^{T}\cdot D_{p}h}(M_{2}^{-1}v^{T})\left(D_{p}h + \frac{x \cdot D_{p}h}{1 - \vert x \vert^{2}}x\right)\right] \\
      & = \frac{1}{h}\left[I + \frac{x\otimes x}{1 - \vert x \vert^{2}}\right.  \\
      & \quad \left. + \frac{1}{1 - (M_{2}^{-1}v^{T})^{T}\cdot D_{p}h}\left(h^{-1}\left(\frac{\rho}{1 - \vert x \vert^{2}}x^{T} - (D\rho)^{T}\right)\left(D_{p}h + \frac{x \cdot D_{p}h}{1 - \vert x \vert^{2}}x\right)\right)\right],
  \end{align*}
  where we have used the fact that for the row vectors $a,b,c,d$, $(a \otimes b)(c \otimes d) = (b \cdot c)(a \otimes d)$. Denoting that
  \begin{equation*}
    A = \frac{1}{1 - (M_{2}^{-1}v^{T})^{T}\cdot D_{p}h}\left(h^{-1}\left(\frac{\rho}{1 - \vert x \vert^{2}}x^{T} - (D\rho)^{T}\right)\left(D_{p}h + \frac{x \cdot D_{p}h}{1 - \vert x \vert^{2}}x\right)\right),
  \end{equation*}
  then from
  \begin{equation*}
    (M_{2}^{-1}v^{T})^{T}\cdot D_{p}h = h^{-1}\left[\frac{\rho}{1 - \vert x \vert^{2}}(x \cdot D_{p}h) - D\rho \cdot D_{p}h\right],
  \end{equation*}
  we have
  \begin{equation*}
    A = \frac{1}{h - \left[\dfrac{\rho}{1 - \vert x \vert^{2}}(x \cdot D_{p}h) - D\rho \cdot D_{p}h\right]}\left(\frac{\rho}{1 - \vert x \vert^{2}}x - D\rho\right)\otimes \left(D_{p}h + \frac{x \cdot D_{p}h}{1 - \vert x \vert^{2}}x\right),
  \end{equation*}
  so (\ref{5.25}) can be written as
  \begin{equation*}
    C^{-1} = \frac{1}{h}\left[I + \frac{x \otimes x}{1 - \vert x \vert^{2}} + A\right].
  \end{equation*}

  From (\ref{5.22}), we have
  \begin{equation*}
    Dy = \frac{1}{\kappa}C(C^{-1}B + D^{2}\rho),
  \end{equation*}
  hence
  \begin{equation}\label{5.29}
    \det Dy = \frac{1}{\kappa^{n-1}} \det C \det (C^{-1}B + D^{2}\rho).
  \end{equation}
  Combining (\ref{5.2}), (\ref{5.5}), (\ref{5.17}) and (\ref{5.25}), we have
  \begin{equation*}
    \vert \frac{1}{\kappa^{n-1}} \det C \det (C^{-1}B + D^{2}\rho) \vert = \frac{f(x)t_{\Gamma}(x)\omega}{\kappa g(T(x))},
  \end{equation*}
  then from (\ref{5.28}), we finally obtain
  \begin{equation*}
    \vert \det (D^{2}\rho + C^{-1}B) \vert \leq \frac{f(x)t_{\Gamma}(x)\vert \kappa \vert^{n - 2} \omega}{g(T(x))h^{n - 1}\left(1 - h^{-1}\left(\dfrac{\rho}{1 - \vert x \vert^{2}}x - D\rho\right)\cdot D_{p}h\right)}.
  \end{equation*}

\end{proof}

\end{document}
